\subsection{CARTAN SUBALGEBRAS OF COMPACT ALGEBRAS}

Lemma 1. Let G be a Lie group, K a compact subgroup of G, and F an invariant bilinear form on L(G). Let \(x, y \in \mathrm{L}(\mathrm{G})\) . There exists an element k of K such that \(\mathrm{F}(u, [(\mathrm{Ad}k)(x), y]) = 0\) for all \(u \in \mathrm{L}(\mathrm{K})\) .

The function \(v \mapsto \mathrm{F}((\operatorname{Ad} v)(x), y)\) from K to R is continuous, and hence has a minimum at some point \(k \in K\) . Let \(u \in \mathrm{L}(K)\) and put

\[
h (t) = \mathrm{F} ((\mathrm{Ad} \exp (t u). k) (x), y), \quad t \in \mathbf {R}.
\]

We have \(h(t) \geq h(0)\) for all \(t\) ; moreover, by Chap. III, §3, no. 12, Prop. 44,

\[
\frac {d h}{d t} (0) = \mathrm{F} ([ u, (\mathrm{Ad} k) (x) ], y) = \mathrm{F} (u, [ (\mathrm{Ad} k) (x), y ]),
\]

hence the lemma (Functions of a Real Variable, Chap. I, §1, no. 7, Prop. 7).

THEOREM 1. Let g be a compact Lie algebra. The Cartan subalgebras of g (Chap. VII, §2, no. 1, Def. 1) are its maximal commutative subalgebras; in particular, g is the union of its Cartan subalgebras. The group \(\operatorname{Int}(\mathfrak{g})\) operates transitively on the set of Cartan subalgebras of g.

Since g is reductive, its Cartan subalgebras are commutative (Chap. VII, §2, no. 4, Cor. 3 of Th. 2). Conversely, let t be a commutative subalgebra of g. By §1, no. 3, Prop. 1, ad x is semi-simple for all \(x \in t\) ; by Chap. VII, §2, no. 3, Prop. 10, there exists a Cartan subalgebra of g containing t. This proves the first assertion of the theorem.

Now let \(\mathfrak{t}\) and \(\mathfrak{t}'\) be two Cartan subalgebras of \(\mathfrak{g}\) . We prove that there exists \(u \in \operatorname{Int}(\mathfrak{g})\) such that \(u(\mathfrak{t}) = \mathfrak{t}'\) . By Prop. 1 of §1, no. 3, we can assume that \(\mathfrak{g}\) is of the form L(G), where G is a connected compact Lie group, and can choose a separating invariant symmetric bilinear form F on g. Let x (resp. \(x'\) ) be a regular element of g such that \(\mathfrak{t} = \mathfrak{g}^{0}(x)\) (resp. \(\mathfrak{t}' = \mathfrak{g}^{0}(x')\) ) (Chap. VII, §3, no. 3, Th. 2). Applying Lemma 1 with K = G, we see that there exists \(k \in G\) such that \([(\operatorname{Ad} k)(x), x']\) is orthogonal to g with respect to F, and hence is zero; then \((\operatorname{Ad} k)(x) \in \mathfrak{g}^{0}(x') = \mathfrak{t}'\) , so \(\mathfrak{g}^{0}((\operatorname{Ad} k)(x)) = \mathfrak{t}'\) since \((\operatorname{Ad} k)(x)\) is regular. We conclude that \((\operatorname{Ad} k)(\mathfrak{t}) = \mathfrak{t}'\) , hence the theorem.

COROLLARY. Let t and \(t'\) be Cartan subalgebras of g, a a subset of t, and u an automorphism of g that takes a into \(t'\) . There exists an element v of \(\text{Int}(\mathfrak{g})\) such that \(u \circ v\) takes t to \(t'\) , and coincides with u on a.

Put \(\mathrm{G} = \operatorname{Int}(\mathfrak{g})\) , and consider the fixer \(Z_{\mathrm{G}}(\mathfrak{a})\) of \(\mathfrak{a}\) in \(G\) ; this is a Lie subgroup of \(G\) , whose Lie algebra \(\mathfrak{z}_{\mathfrak{g}}(\mathfrak{a})\) consists of the elements of \(\mathfrak{g}\) that commute with every element of \(\mathfrak{a}\) (Chap. III, §9, no. 3, Prop. 7). Then \(t\) and \(u^{-1}(t')\) are two Cartan subalgebras of the compact Lie algebra \(\mathfrak{z}_{\mathfrak{g}}(\mathfrak{a})\) . By Th. 1, there exists an element \(v\) of \(Z_{\mathrm{G}}(\mathfrak{a})\) such that \(v(t) = u^{-1}(t')\) ; any such element has the desired properties.

